\clearpage
\section*{The City Council Journals complete the zoning record before November 2010}
\textit{September 28, 2026. Agent-drafted work on branch \texttt{explore-alderman-measures}; Jacob asked to go back
before 2010, starting with a 2008 pilot and eventually to 2000, and to OCR the Journals afresh.}

eLMS records no zoning map amendment introduced before April 2010 and only 1 to 5 a month from April to October
2010, while the Journals of the Proceedings list 11 to 57 a month. The stall and timing measures therefore rest on
applications from November 2010 on, and the Journals are the only complete source before then. Each Journal lists
the amendments introduced at a meeting, by applicant or alderman, with the map sheet, the districts before and after,
the boundary and the common address, and prints each ordinance passed with its application number and change.
Journals for 2008--2011 are downloaded and parsed into 1,339 introductions and 1,179 passages.

The PDFs' text layers are the City Clerk's own OCR, which misreads district codes and words (``Cl-2,'' ``82-3,''
``RMS,'' ``Applicafion''), so the 2,635 zoning pages of 2008--2011 were read again with tesseract. Tesseract makes
fewer errors but different ones (``ail'' for ``all,'' ``|'' for the map letter I, ``Reciassification''), and the
parser allows for both. On the 87 passages of 2008 where the two texts' parses disagreed, read from the page images,
the fresh OCR's parse matched the printed page on 80 map sheets, 81 districts before and 83 districts after, against
72, 67 and 63 for the text layer's. The parser stops if it contradicts any of 28 entries read from the printed pages.

Where both sources exist, from November 2010 through 2011, the Journals and eLMS agree closely: 372 of the Journals'
380 introductions match an eLMS amendment by record number, and 372 of eLMS's 399 amendments are found in the
Journals. The application numbers agree for 248 of 249 matched introductions and the map sheets for 331 of 335; the
districts before and after agree for about 89 percent, and most disagreements are differences between the sources
rather than misreadings: a planned development recorded separately in eLMS, or a request changed between its
introduction and the ordinance.

The introductions print the Clerk's record number from December 17, 2008 and the application number from April 22,
2009, so from mid-2009 each introduction can be followed to its passage by number. Earlier introductions must be
matched to passages by map sheet and boundary, with hand review, which is the next step, followed by outcomes for
the 2008 applications through the end of that term in May 2011. Going back to 2000 will need the district codes of
the zoning ordinance before its 2004 rewrite.

\textit{Update, September 28, 2026: four checks} (\path{tasks/audits/zoning_record_validation}). The gap is in
eLMS itself, not in the download: eLMS holds 0 to 74 matters of any kind a month before November 2010 and none of the
173 record numbers the Journals print for January--October 2010. On 32 randomly drawn pages of 2008 and 2009, read
by hand, the parser finds all 54 introductions and 30 passages and no others, and reads the map and districts of
53--54 introductions and 26--29 passages correctly. Against eLMS's passages of November 2010 through 2011 the
Journal's dates, maps and districts agree for 263 of 267, 236 of 237, 170 of 175 and 179 of 191 matched ordinances,
after correcting three Journals the City Clerk's listing dates wrongly. The main gap the checks found: an ordinance
the Committee on Zoning reports on its own is printed under the committee's heading and not read, 10 to 15 a year
from 2009, many of them planned developments; those reports also carry withdrawals. Completing the pre-2010 dataset
also needs introductions linked to outcomes, wards, direction and refilings; it cannot have project-size controls,
since the application forms exist only in eLMS.

\textit{Update, later on September 28, 2026: the Council's actions.} Reading those committee reports showed that the
parser had taken every printed ordinance for a passage, though a report can also defer an ordinance (to be taken up
at a later meeting), withdraw it or re-refer it. The parser now reads the Committee on Zoning's reports themselves,
including those it makes on a single ordinance, and records the Council's action on each ordinance and the
withdrawals and deferrals a report notes without printing the ordinance. In 2008--2011 it finds 1,229 ordinances:
1,210 passed, 5 deferred and 14 withdrawn. Against eLMS from November 2010 through 2011, 280 passed ordinances match
an eLMS passage, with dates agreeing for 276 of them, and eLMS records all 7 matched withdrawals as withdrawn. The
City Clerk's listing had dated three Journals wrongly; each Journal's date is now read from its own pages, and the
introductions' dates agree with eLMS for all 372 matched. Linking each introduction to its outcome is the next step:
applications carry their application numbers from April 2009, but the 355 earlier applications and all 432
aldermen's amendments of 2008--2011 must be matched on map sheet, districts and boundary, with hand review.

\textit{Update, evening of September 28, 2026: introductions linked to their outcomes}
(\path{tasks/link_journal_zoning_outcomes}). Each ordinance is linked to the introduction it resolves by record number
(from November 2010), by application number (from April 2009), by the words its boundary shares with an earlier
introduction's, and, for an application still unlinked, by its number's place in the numbering: applications are
numbered in order and each meeting's take the next run, so a number places an application between two meetings. The
linking uses no hand decisions. It is checked three ways. With every printed number hidden, the boundary and
numbering rules link 421 of the 433 applications whose introductions the numbers identify, all correctly. The numbers
of all 696 linked applications, including 255 of 2008 and early 2009 linked by boundary alone, fall between those of
the meetings before and after their introductions'. And for the 377 introductions of November 2010--2011 that eLMS
also holds, the outcome as of December 2011 agrees for 375; one of the two disagreements is a wrong link, an
alderman's amendment introduced directly in committee that is linked to an earlier introduction of the same site.

Reading the unlinked ordinances found parser faults, since fixed: an ordinance heading without ``On,'' a curly quote
inside ``to classify as,'' record numbers after a stray mark, and an aldermen's section whose opening has a stray
period. The Journals now match every eLMS passage of November 2010--2011. It also found two facts that bear on the
measures. Matters did not lapse at the term change of May 2011: 48 amendments introduced from January to May 2011 and
5 from 2009 and 2010 passed after it, so ``stalled'' here means not decided at least a year after introduction,
within the Journals read (through December 2011). And aldermen introduce some amendments directly in committee, to
pass at the same meeting, so those ordinances have no introduction.

Applications stalled far more often before 2011 than after: 23 percent of those introduced in 2008, 21 percent in
2009 and 14 percent in 2010, against 6 to 8 percent a year in eLMS from 2010 to 2015. A year of follow-up overstates
stalls: of 19 applications of November 2010 to May 2011 undecided in December 2011, 5 passed in 2012. That accounts
for much of the 2010 figure but little of 2008 and 2009, which have two to four years of follow-up. Reading the
ordinances of the 2012 and 2013 Journals would give every introduction of 2008--2010 at least two years of follow-up.
The Journals of 1981--2011 are now downloaded (698 files; the Clerk's listing shows 25 a page, which the first
download missed).

\textit{Update, night of September 28, 2026: wards from boundaries}
(\path{tasks/place_journal_zoning_amendments}). The Journals print no address with an amendment before July 9, 2008,
and aldermen's amendments rarely one after, only the map sheet and the boundary. So every amendment is placed from its
boundary, the same way in every year: the streets it names are matched to the City's 2013 street centerlines, the
point is set at their corners, and, where a parcel has only one or two corners, moved to the side of each street
that the boundary gives (``a line 131.20 feet north of \ldots\ West Cermak Road''), since streets are often ward
boundaries. Of 1,341 introductions of 2008--2011, 1,246 are placed; 72 of 906 applications are not, mostly
boundaries described by metes and bounds or with streets OCR misreads.

The City's zoning map turned out to record, for each district set since about 2002, the application number and
passage date of the amendment that set it (\path{tasks/download_chicago_gis_layers}, the map in force and that of
November 2012). A district stays on the map only until amended again, and an application that stalled never reaches
it, so the map can check the placement but not replace it. For 525 passed ordinances matched to districts by number,
date and district class, the placed point is inside the parcel for 371 and within 237 feet for 99 percent, and the
placed ward is the parcel's for all 449 parcels in one ward (\path{tasks/audits/zoning_record_validation}). The placed
ward is also the filing ward of 402 of 411 aldermen's amendments. The map sometimes files a district under another
amendment's number, so a match must also agree on the district's class.

\textit{Update, September 29, 2026: the Journals of 2000--2007.} The parser, links and places now cover 2000--2011:
6,924 introductions and 5,806 ordinances. Before November 1, 2004, the Journals rezone under the ordinance of 1957,
whose codes overlap the 2004 ordinance's with different districts (B3-2 was General Retail and is Community Shopping);
its codes are kept apart, marked ``1957:'', by date and by the district name printed with them. The degraded scans
of 2002 print the section headings with stray marks, and letter-spaced text layers of 2001 hid zoning pages from the
page selection; both are now read. A random sample of 64 pages of 2000--2007, read blind from the images, has 111
introductions and 55 ordinances, and the parser finds every one and no others; it reads the districts of all but one
(an illegible ``R\textcent'') and the map of all but one (an ambiguous ``5-]''). The first scoring of the sample found
11 general faults, since fixed: a period after a lettered street (``South Avenue M.'') taken for an initial's, maps
printed ``Map Number11-L'', articles run into codes (``aC1-1''), and four wordings of an ordinance's change. OCR also
reads ``R5 General Residence'' as ``RS'' 533 times; it is repaired.

With the introductions of 2007 read, the ordinances of early 2008 now link: 486 of 5,806 ordinances are unlinked,
154 of them passed in 2000 on introductions of 1999, and 194 of the rest aldermen's amendments, many introduced
directly in committee. Stall rates stay high: 14 to 23 percent of applications introduced in each year of 2000--2010
and 26 to 46 percent of aldermen's amendments were not decided by December 2011. The applications that remain unlinked
could lower the applications' rate by at most about 4 points a year, so the gap with eLMS's 6 to 8 percent is not a
linking artifact. Where only boundaries link them (before 2009), a site introduced again after an earlier
introduction stalled can be linked to the earlier one; among the 515 ordinances linked by number, the boundary rule
does so for 3.

Before May 5, 2003, amendments are placed in the wards of 1998. The zoning map matches 1,684 passed ordinances, and
the placed ward is the parcel's for 1,444 of the 1,448 parcels in one ward. The filing wards of aldermen's
amendments revealed a transition: the Council redrew the wards on December 19, 2001, for the term that began in May
2003, and from May 2002 aldermen filed amendments in the wards they would represent. Where the two maps differ, the
filing ward is the 1998 map's for 30 of 33 amendments from September 2001 to March 2002, and the 2003 map's for 36 of
46 from May 2002 to April 2003. The ward and alderman are assigned by the map in force; whether an application of
that year belongs to the sitting or the incoming alderman is a choice for the analysis, and the same question arises
for the wards redrawn in 2012 and in force from 2015.

\textit{Update, later on September 29, 2026: follow-up and refilings} (\path{tasks/follow_journal_zoning_amendments}).
Two things could make the Journals' stall rates too high: amendments undecided in December 2011 that passed later,
and applications that stalled but were filed again and passed. eLMS follows 103 of the 1,643 introductions undecided
in December 2011 by record or application number; 68 later passed. With that follow-up, 18 of the Journals' 215
applications of 2011 never passed, 8.4 percent, which is eLMS's own rate for 2011: where the two sources overlap,
the Journals measure stalls as eLMS does. A stalled amendment is refiled by the first later amendment of the same
site by the same kind of filer within a council term: a nearly identical boundary, a similar boundary placed within
300 feet, the same application number, or, from July 2008, an overlapping common address. Each of 191 candidate pairs
was read by hand (two readers, disagreements adjudicated). 150 of 1,594 amendments that never passed were refiled,
110 of them successfully. Net of those, 14 to 22 percent of applications introduced in each year of 2000--2008 never
passed, against 7 percent in 2011 and 4 to 8 percent a year in eLMS. Before July 2008 refilings of large sites can be
missed, since only boundaries and placed points identify them; if stalled applications were refiled and passed as
often as in eLMS (24 percent), the rates would still be 11 to 18 percent. That stalls were about twice as common in
the 2000s as since 2011 appears to be a real change rather than a feature of the source.

The follow-up also found that the Census geocoder, used for eLMS title addresses
(\path{tasks/assign_zoning_amendment_wards}), sometimes returns the house number on the opposite side of the city
(5145 N California Ave as 5145 S California Ave) or in Chicago Heights. Such matches are now rejected and the address
placed on the City's street centerlines instead; 25 were rejected, and 22 amendments changed ward and alderman. The
alderman zoning measures moved little (shrunk stall effects correlate at 0.998 with the earlier ones).

\textit{Update: placing more amendments.} 7 percent of introductions had no place, many of them ordinary corner
lots whose boundary names a street by a type it does not have at that corner: Diversey Avenue is Diversey Parkway at
Lakewood, Fulton Street is Fulton Boulevard at Sacramento, and the Journals misprint types (``South Sangamon Avenue'').
An amendment not placed with the printed types is now tried again with the through-street types taken as one, with
drives, places and courts kept apart (South Peoria Drive is not Peoria Street). 131 more introductions are placed,
leaving 5 percent unplaced: boundaries naming fewer than two streets (metes and bounds, railroad lines), and streets that do not meet as printed, some with a direction misprinted (``North Halsted Street'' at West 33rd). All 35 of the new placements the zoning map can check lie within 240 feet of their parcels, and their wards
agree for all 31 in one ward.

\textit{Update: missing boundaries, addresses and directions.} The introductions left unplaced were not a random
5 percent: 44 percent of the applications naming fewer than two streets stalled, against 17 percent of those placed,
and more were planned developments, so leaving them out would favor aldermen whose wards receive large irregular
sites. Three general repairs followed. The parser now reads a boundary printed after ``bounded'' without ``by:'', ``bound
by'', ``bounded \copyright\ by:'' or, failing those, ``as follows:''; 58 introductions had none, 21 still do. With those
boundaries, 26 more ordinances link to their introductions, and 27 applications counted as stalled in fact passed,
lowering the stall rates of 2000--2007 by up to a point and a half a year. The placement now uses an introduction's
common address (from July 2008) when its streets fail, and tries each other direction for one street when the
printed streets do not meet, accepting only a unique repair. 3.4 percent of introductions remain unplaced (mostly lot
descriptions and metes and bounds) and 4 percent of applications have no ward; all 37 of the new placements the
zoning map can check lie within 330 feet of their parcels. Net of refilings that passed, 14 to 22 percent of
applications introduced in 2000--2008 never passed.

\textit{Update, night of September 29, 2026: a blind search for stalled applications' passage}
(\path{tasks/audits/zoning_record_validation}, section 7 of its README). Whether an application recorded as stalled in
fact passed was tested without the parser or the links. 40 applications of 2000--2008 recorded as stalled were drawn
at random, with 10 recorded as passed as controls. Every later page of the Journals through 2011 was searched in its
own text for each application's distances, streets, map sheet and applicant's name, and the four best-matching pages
were read from their images by readers who did not know which applications were controls. They found the passage of 8
of the 10 controls (the other two were an ordinance amended to name another street and a page deep in a planned
development's exhibit). Of the 40 stalled, 31 had no later passage and 2 passed through their recorded refilings; 5
had passed and were wrongly recorded, for three reasons, each now fixed in general. OCR had printed ``area-bounded
by'' or ``bounded. by'', so the parser read no boundary to link (\path{tasks/parse_journal_zoning_amendments});
tesseract had split a page into two columns and broken its ordinance heading, and 9 such pages are now read again as
one column (\path{tasks/ocr_council_journal_zoning_pages}); and an application printed twice at one meeting, as the
applicant's and as the sponsoring alderman's, had its ordinance linked to only one printing, and now shares it with
its twin (\path{tasks/link_journal_zoning_outcomes}). Together these turned 32 introductions of 2000--2011 from
stalled to passed. Net of refilings that passed, 13 to 22 percent of applications introduced in each year of
2000--2008 never passed, against 7 percent in 2011; with eLMS's rate of refiling, 10 to 17 percent. The five errors
were found in this sample and fixed from it, so it cannot measure the error that remains; a fresh draw would.

\textit{Reproduction.} Run \texttt{make} in \path{tasks/link_journal_zoning_outcomes/code/} and
\path{tasks/place_journal_zoning_amendments/code/}, which read \path{tasks/parse_journal_zoning_amendments}, and that in
turn \path{tasks/ocr_council_journal_zoning_pages} and \path{tasks/extract_council_journal_text}; the task READMEs give
the counts by year and the validation.
